\defmodule{MRG}

This module offers tools to analyze simple or combined 
multiple recursive (linear congruential) generators (MRGs).
One must first initialize the module with a given modulus,
a given order, and a maximal dimension, by calling {\tt InitMRG}.
Each MRG will then be defined by a vector of multipliers {\tt a},
where {\tt a[i]} represents $a_i$ and the MRG is defined by
 $$ x_n := (a_1 x_{n-1} + \cdots + a_k x_{n-k}) \mod m.$$
One can then build the lattice basis associated to a vector of multipliers
for a given dimension, or apply the spectral test, or the Beyer test,
or both simultaneously, for all dimensions up to a specified limit.
The multipliers and modulus can be represented as {\tt LONGINT} or
{\tt SuperInteger}.  This can be changed by the command {\tt select}.

\bigskip\hrule
\code

DEFINITION MODULE MRG;

FROM MULT     IMPORT MScal, MVect;
FROM BASIS    IMPORT BVect;
FROM CONFIG   IMPORT MDim, BDim, OutputType;
FROM LATBASIS IMPORT Basis;
FROM TESTLAT  IMPORT TestConfig;
FROM NORM     IMPORT CriterionType, NormType, Normaliz;
FROM REDBAS   IMPORT VerifyType;
FROM SUPINT   IMPORT SuperInteger;
FROM SUPFACT  IMPORT Factors;
IMPORT BASISLR, MULTLI, MULTSI;

TYPE
   MRGen =  POINTER TO InfoMRG;
   InfoMRG = RECORD
      Bas   : Basis;
      k     : MDim;                     (* k = Order  of the recurrence *)
      aLI   : MULTLI.MVect;             (* Coefficients of the recurrence *)
      aSI   : MULTSI.MVect;             (*               "                *)
      END;



   LatticeType = (Full, Recurrent, Orbit, PrimePower);
\endcode
 \tabb Indicates which type of lattice is analyzed: that generated 
  by all possible states ({\tt Full}), 
  or by the set of recurrent states ({\tt Recurrent}), 
  or by a single forward orbit of the generator
  ({\tt Orbit}) for the case where it has several cycles, 
  or by a single orbit in the case where
  some component is an MLCG with modulus that is a power of a prime $p$
  ({\tt PrimePower}).  In the latter case, one must divide the modulus
  by an appropriate power of $p$, as described in \cite{rLEC97c},
  before calling a test procedure. 
  When {\tt LatticeType = Orbit}, the initial state must be specified
  in the array {\tt InSta} (see below).
  See also the data field {\em LatticeType\/} for the program {\tt seekl}.
 \endtabb
\code


VAR
   Bas : Basis;
\endcode
 \tabb  Lattice basis built and used by the procedure {\tt Test}.
   Can be recovered after calling {\tt Test}.
 \endtabb
\code

   Q2, D2, S2 : BASISLR.BVect;
\endcode
 \tabb  Those vectors contain the {\em squares\/} of the Beyer 
   quotients $q_t$, of the distances between
   hyperplanes $d_t$, and of the values of $S_t$, respectively, computed in
   the last call of {\tt Test} or {\tt TestLa}.
 \endtabb
\code

   Mins2, Minq2 : BASISLR.BVect;
\endcode
 \tabb  Gives the current minimal values of 
   $S_t$ and $q_t$ for a generator to be considered, in each dimension.
   The module initializes these bounds to zero.
 \endtabb
\code
 
   MaxAllDim : BOOLEAN;
\endcode
 \tabb  {\tt TRUE} if we are always looking for generators with the
  highest figures of merit, in all the dimensions.
  {\tt FALSE} if we want to {\em minimize\/} the figure of merit in
  some set of dimensions (that is, in some interval or category).
 \endtabb
\code

   TestCompleted : BOOLEAN;
\endcode
 \tabb  {\tt TRUE} if the last call to {\tt Test} or {\tt TestLa}
  has been successful, {\tt FALSE} otherwise.
 \endtabb
\code

   Lacunary : BOOLEAN;
\endcode
 \tabb  When this variable is {\tt TRUE}, the module treats lacunary
   indices; otherwise it assumes that the vectors are formed with
   successive indices.
 \endtabb
\code

   LatType : LatticeType;
\endcode
 \tabb  Indicates which lattice or sublattice is analyzed.
  See the data field {\em LatticeType\/} in {\tt seekl}.
  However, in the case {\tt PrimePower}, the modulus of each component
  must be divided appropriately {\em before\/} calling {\tt Test}
  or {\tt TestLa}; these test procedures assume that the modulus has 
  already been ajusted (reduced).
 \endtabb
\code

   DualLattice : BOOLEAN;
\endcode
 \tabb When {\tt TRUE}, we seek short vectors in the {\em dual\/}
  lattice (e.g., to compute the distance between adjacent hyperplanes).
  When {\tt FALSE}, we seek short vectors in the {\em primal\/}
  lattice.  The default value is {\tt TRUE}.
 \endtabb
\code

   UsedNorm : NormType;
\endcode
 \tabb  Which norm type is used to measure the length of the vectors.
  The default value is {\tt L2Norm}.
  WARNING: For other norms, the implementation is only partial, and
  only {\tt Spectral} is admitted for the type of test.
 \endtabb
\code

  NormalG : Normaliz;
\endcode
 \tabb To normalize the length of the shortest vectors (See module NORM).
 \endtabb
\code

   InvertCrit : BOOLEAN;
\endcode
 \tabb  If {\tt TRUE}, the value of the criterion $S_t$ is {\em inverted\/} 
  in the results, i.e., the criterion $S_t$ becomes a number larger than 1,
  the {\em smaller\/} the better.
 \endtabb
\code

   InSta : ARRAY MDim OF SuperInteger;    
\endcode
 \tabb If {\tt LatType = Orbit} then the user must specify a seed which
  will be stored in {\tt InSta}. 
 \endtabb
\code

PROCEDURE InitMRG (m : MScal;  k : MDim;  d : BDim);
\endcode
  \tab  From now on, this module will treat multiple recursive generators
   of order {\tt k} with modulus {\tt m}. 
   Dimensions of bases cannot exceed {\tt d}.
   (One can recall this procedure to change {\tt m}, {\tt k}, or {\tt d}).
   This procedure calls {\tt MULT.MSetm}.
\endtab
\code

PROCEDURE DeleteMRG (d : BDim);
\endcode
  \tab Deallocates the memory allocated to the matrices of {\tt SuperIntegers}
  {\tt VSI} and {\tt WSI}, and to the vectors of type {\tt BVect}, 
  {\tt xi, LacI}
  and {\tt q}, at a previous call to {\tt InitMRG}. It also calls 
  the procedure {\tt LATBASIS.DeleteBasis}, because {\tt LATBASIS.InitBasis} is
  called by {\tt InitMRG}. Must be called only if {\tt InitMRG} has been called
  previously.
\endtab
\code

PROCEDURE BuildBasis (B : Basis;   VAR a : MVect;   n : BDim);
\endcode
   \tab  Builds the basis directly in dimension {\tt n}.
   \endtab
\code

PROCEDURE IncDimBasis (B : Basis);
\endcode
   \tab  Increments by 1 the dimension of the basis.
   \endtab
\code

PROCEDURE Test
   ( VAR a          : MVect;
     T              : CriterionType;
     FromDim, ToDim : BDim;
     VerifyBB       : VerifyType;
     Out            : OutputType );

\endcode
 \tab  Applies the test {\tt T} (Beyer and/or Spectral)
   to the MRG generator corresponding to {\tt a}, in dimensions
   {\tt FromDim} to {\tt ToDim}.
   The initial basis is built automatically, and during the procedure,
   the variable {\tt Bas} will contain the current basis.
   {\tt VerifyBB} indicates if
   the results must be verified in a second pass, with error bounds
   on the floating-point computations for the branch-and-bound.
   The parameter {\tt Out} specifies the form of output produced.
   If a test fails, then the procedure stops and the variable
   {\tt TestCompleted} is set to {\tt FALSE}.
 \endtab
\code

PROCEDURE TestCri
   ( VAR a          : MVect;
     T              : CriterionType;
     FromDim, ToDim : BDim;
     VerifyBB       : VerifyType;
     Out            : OutputType;
     r              : BDim;
     n              : INTEGER );
\endcode
  \tab Does the same thing as {\tt Test}, but two additional parameters
  {\tt r} and {\tt n} must be specified, which are the rank and the copy 
  factor for copy-rules, and then the basis are modified accordingly.
  Must be used only if {\tt FromDim} is larger than {\tt r}. Not very useful
  in the context of LCG's, but it would have been complicated to write
  this procedure and the one it calls, {\tt CriSpectralTest()}, in a
  separate module, because they use variables internally declared in 
  {\tt MRG}. The LCG corresponding to the rank 1 rule must have its
  modulo and multiplier modified before {\tt TestCri} is called.
  THIS PROCEDURE IS NOT IMPLEMENTED. 
  \endtab
\code

PROCEDURE BuildLaBasis
   (B : Basis;   VAR a : MVect;   VAR I : BVect;   n, MaxDim : BDim);
\endcode
 \tab  Builds the (lacunary) basis directly in dimension {\tt n}.
   One must indicate MaxDim and use IncDimLaBasis only for dimensions
   smaller than this.
   The vector {\tt I} must contain the {\tt MaxDim} lacunary indices.
 \endtab
\code

PROCEDURE IncDimLaBasis (B : Basis);
\endcode
  \tab  Increments by 1 the dimension of the basis. The basis
   must have been initialized by {\tt InitLaBasis}
   and the dimension of {\tt B} must be less than {\tt MaxDim}.
  \endtab
\code

PROCEDURE TestLa
   ( VAR a          : MVect;
     VAR I          : BVect;
     T              : CriterionType;
     FromDim, ToDim : BDim;
     VerifyBB       : VerifyType;
     Out            : OutputType );
\endcode
 \tab  Similar as {\tt Test}, but with lacunary indices.
  The vector {\tt I} must contain the {\tt ToDim} lacunary indices.
 \endtab
\code

PROCEDURE PerMax (VAR a : MVect;   VAR ListmMS1, Listr : Factors) : BOOLEAN;
\endcode
  \tab  Checks if the generator corresponding to {\tt a}
    has maximal period. If {\tt m} is not prime then the procedures stops
    and {\tt TestCompleted} is set to {\tt FALSE}.
    {\tt ListmMS1} and {\tt Listr} must contain the factors of
    {\tt m-1} and {\tt r}, or {\tt NIL} if the factors are unknown.
    (See {\tt MULT.SetFactors}.)
   \endtab
\code

END MRG.
\endcode

